%6.1 iteration을 (\theta^{(k+1)} = T(\theta^{(k)})) 형태의 연산자로 정리
%6.2 (T)의 Lipschitz bound를 (H_\phi, P_\psi)의 Lipschitz 상수로 상계
%6.3 spectral normalization이 이를 어떻게 보장하는지(정리/조건)
%6.4 Theorem: unique fixed point + exponential convergence
%6.5 Corollary: K=10 고정이 충분한 이유(오차 bound로 연결 가능)

\section{Theory 1: Convergence}
\label{sec:convergence}

In this section we analyze the iteration operator
$\theta^{(k+1)} = \mathcal{T}(\theta^{(k)};\mathbf{x}_\mathrm{obs})$
introduced in \ref{sec:method}.
The introduction of the iteration raises a critical theoretical questions:

\emph{Convergence: Does this iteration process stably converge to a unique fixed point?}
\emph{Sensitivity: Is it sensitive to the initial guess $\theta^{(0)}$?}

In this section, we prove that by enforcing strict Lipschitz bounds on two neural networks $H_\phi$ and $P_\psi$, the composite operator $\mathcal{T}$ becomes a global contraction mapping, guaranteeing exponential convergence to a unique fixed point $\theta^*$ from any arbitrary initialization.

\subsection{Preliminaries: contraction mappings and Banach's theorem}
\label{subsec:banach}

Let $(\Theta,\|\cdot\|_2)$ be a complete metric space.
\footnote{In practice, $\Theta$ is taken as a closed subset of $\mathbb{R}^p$ in normalized coordinates.}
A map $\mathcal{T}:\Theta\to\Theta$ is a \textit{contraction} if there exists $K\in[0,1)$ such that

\begin{equation}
\label{eq:contraction-def}
    \lVert \mathcal{T}(\theta_a)-\mathcal{T}(\theta_b)\|_2 \le K\|\theta_a-\theta_b\rVert_2,
    \qquad \forall\,\theta_a,\theta_b\in\Theta.
\end{equation}

By \emph{Banach's fixed-point theorem}, a contraction has a unique fixed point $\theta^*$ and the iteration
$\theta^{(k+1)}=\mathcal{T}(\theta^{(k)})$ converges to $\theta^*$ for any initialization. 
Moreover, the convergence is geometric:

\begin{equation}
\label{eq:geom-rate}
    \|\theta^{(k)}-\theta^*\|_2 \le K^k\|\theta^{(0)}-\theta^*\|_2.
\end{equation}

\subsection{A Sufficient Condition for Contractivity of $\mathcal{T}$}
\label{subsec:contractivity}

Recall that our update operator is the composition:

\begin{equation}
\label{eq:T-repeat}
    \mathcal{T}(\theta;\mathbf{x}_\mathrm{obs}) = P_\psi\big(\mathbf{x}_\mathrm{obs}, H_\phi(\mathbf{x}_\mathrm{obs},\theta)\big).
\end{equation}

We provide a simple sufficient condition for \eqref{eq:contraction-def} in terms of Lipschitz constants of the two networks.

\begin{assumption}[Lipschitz continuity of the learned maps]
\label{ass:lipschitz}
    For each fixed input $\mathbf{x}_\mathrm{obs}$, the hidden-state estimator $H_\phi$ is Lipschitz in $\theta$:
    there exists $\mathrm{Lip}_H\ge0$ such that
    $$
    \lVert H_\phi(\mathbf{x}_\mathrm{obs},\theta_a)-H_\phi(\mathbf{x}_\mathrm{obs},\theta_b)\rVert_2 \le 
    \mathrm{Lip}_H\|\theta_a-\theta_b\|_2,
    \qquad \forall\,\theta_a,\theta_b\in\Theta.
    $$
    Moreover, the parameter estimator $P_\psi$ is Lipschitz in its hidden-state input:
    there exists $\mathrm{Lip}_P\ge0$ such that for all $\mathbf{z}_a,\mathbf{z}_b$ in the range of $H_\phi(\mathbf{x}_\mathrm{obs},\cdot)$,
    $$
    \lVert P_\psi(\mathbf{x}_\mathrm{obs},\mathbf{z}_a)-P_\psi(\mathbf{x}_\mathrm{obs},\mathbf{z}_b)\|_2 \le \mathrm{Lip}_P\|\mathbf{z}_a-\mathbf{z}_b\|_2.
    $$
\end{assumption}

\begin{theorem}[Contractivity of the Iterative Operator]
\label{thm:contraction}
    Under Assumption~\ref{ass:lipschitz}, the update operator $\mathcal{T}(\cdot;\mathbf{x}_\mathrm{obs})$ is Lipschitz with constant at most $\mathrm{Lip}_\mathcal{T}:=\mathrm{Lip}_P\mathrm{Lip}_H$:
    $$
    \|T(\theta_a;\mathbf{x}_\mathrm{obs})-T(\theta_b;\mathbf{x}_\mathrm{obs})\|_2 \le \mathrm{Lip}_P\mathrm{Lip}_H \|\theta_a-\theta_b\|_2,
    \qquad \forall\,\theta_a,\theta_b\in\Theta.
    $$
    In particular, if $\mathrm{Lip}_P\mathrm{Lip}_H<1$, then $\mathcal{T}(\cdot;\mathbf{x}_\mathrm{obs})$ is a contraction on $(\Theta,\|\cdot\|_2)$ and hence has a unique fixed point
    $\theta^*(\mathbf{x}_\mathrm{obs})$.
    Moreover, for any initialization $\theta^{(0)}\in\Theta$, the iteration $\theta^{(k+1)}=\mathcal{T}(\theta^{(k)};\mathbf{x}_\mathrm{obs})$ converges to
    $\theta^*(\mathbf{x}_\mathrm{obs})$ and satisfies the geometric rate bound \eqref{eq:geom-rate} with $\mathrm{Lip}_\mathcal{T}$.
\end{theorem}

\begin{proof}
    Fix $\mathbf{x}_\mathrm{obs}$ and let $\theta_a,\theta_b\in\Theta$.
    By \eqref{eq:T-repeat} and the Lipschitz property of $P_\psi$ in its second argument,
    \begin{align*}
        \lVert \mathcal{T}(\theta_a;\mathbf{x}_\mathrm{obs})-\mathcal{T}(\theta_b;\mathbf{x}_\mathrm{obs})\rVert_2
    &=
    \lVert P_\psi(\mathbf{x}_\mathrm{obs},H_\phi(\mathbf{x}_\mathrm{obs},\theta_a)) - P_\psi(\mathbf{x}_\mathrm{obs},H_\phi(\mathbf{x}_\mathrm{obs},\theta_b))\rVert_2
    \\
    &\le
    \mathrm{Lip}_P\lVert H_\phi(\mathbf{x}_\mathrm{obs},\theta_a)-H_\phi(\mathbf{x}_\mathrm{obs},\theta_b)\rVert_2.
    \end{align*}
    
    Applying the Lipschitz property of $H_\phi$ in $\theta$ yields
    $$
    \|\mathcal{T}(\theta_a;\mathbf{x}_\mathrm{obs})-\mathcal{T}(\theta_b;\mathbf{x}_\mathrm{obs})\|_2 \le \mathrm{Lip}_P\mathrm{Lip}_H\|\theta_a-\theta_b\|_2.
    $$
    Since $\mathrm{Lip}_\mathcal{T}=\mathrm{Lip}_P\mathrm{Lip}_H<1$, then $\mathcal{T}(\cdot;\mathbf{x}_\mathrm{obs})$ is a contraction and the remaining claims follow from Banach's fixed-point theorem.
\end{proof}

\subsection{Enforcing Lipschitz bounds via spectral normalization}
\label{subsec:sn}

Assumption~\ref{ass:lipschitz} can be promoted in practice by controlling the operator norms of the linear layers.
For a linear map $x\mapsto Wx$, its $\ell_2$-Lipschitz constant equals the spectral norm $\|W\|_2$.
Spectral normalization rescales each weight matrix $W$ by an estimate of $\|W\|_2$ so that the resulting layer has a prescribed operator-norm bound,
typically close to~$1$.
This idea was popularized in the deep learning literature as a training-stabilization technique. \cite{miyato2018spectralnormalizationgenerativeadversarial}

For a feed-forward network composed of linear layers and pointwise nonlinearities with Lipschitz constant at most~$1$ (e.g., $\tanh$),
one obtains a simple upper bound on the network's Lipschitz constant by multiplying the spectral norms of the weight matrices.
Consequently, by applying spectral normalization layerwise (and leaving a small margin for strict inequality),
we can ensure $\mathrm{Lip}_H$ and $\mathrm{Lip}_P$ are bounded and achieve $\mathrm{Lip}_P\mathrm{Lip}_H<1$.
Further implementation details are provided in the supplementary material.


Section~\ref{subsec:exp-convergence} verifies this empirically: without spectral normalization the
measured composite Lipschitz constant exceeds one and the iteration fails to contract, whereas across a
range of margins that keep $\mathrm{Lip}_\mathcal{T}<1$ the iteration converges as predicted.

\subsection{Implications for the number of inference steps}
\label{subsec:Ksteps}

Theorem~\ref{thm:contraction} yields a direct justification for using a small, fixed number of inference steps.
If $\mathcal{T}(\cdot;\mathbf{x}_\mathrm{obs})$ is a contraction with constant $K<1$, then after $k$ iterations,
$$
\lVert\theta^{(k)}-\theta^*(\mathbf{x}_\mathrm{obs})\rVert_2 \le K^k\lVert\theta^{(0)}-\theta^*(\mathbf{x}_\mathrm{obs})\rVert_2,
$$
so the error decreases exponentially in $k$.
This explains why in practice a small number of iterations (e.g., $K=10$) suffices, and why often $1$--$2$ iterations already yield a stable estimate.
